%
% Anwendungen der Residuenrechnung
%
\subsubsection{Anwendungen der Residuenrechnung}

%
% hauptwert: Der Hauptwert eines divergenten Integrals
%
Der cauchysche Hauptwert eines divergenten Integrals erscheint auf den
ersten Blick als eine willkürliche Konstruktion, mit der nicht konvergente
Integrale durch einen Trick konvergent gemacht werden.
\emph{Damien Flury}
\index{Damien Flury}%
\index{Flury, Damien}%
zeigt jedoch in Kapitel~\ref{chapter:hauptwert}, dass die Konstruktion
durch die Methoden der Residuenrechnung gut motiviert sind.

%
% kepler: Das dritte keplersche Gesetz
%
Das dritte keplersche Gesetz stellt einen Zusammenhang zwischen Umlaufzeit
und Bahnhalbachse eines Planeten her.
Die Berechnung der Umlaufzeit führt auf ein schwierig zu behandelndes
reelles Integral.
Im Komplexen wird daraus aber eine einfache Rechnung mit Residuen.
\emph{Jan Klarer}
\index{Jan Klarer}%
\index{Klarer, Jan}%
und
\emph{Manuel Kuhn}
\index{Manuel Kuhn}%
\index{Kuhn, Manuel}%
leiten in Kapitel~\ref{chapter:kepler} das dritte keplersche Gesetz auf
diese Weise her.

%
% laplace: Die Laplace-Inversion.
%
Die Laplace-Transformation ist ein besonders nützliches Werkzeug, mit dem
Ingenieure gewöhnliche Differentialgleichungen lösen können.
Der schwierigste Schritt darin ist die Umkehrung, die in der Praxis oft
Probleme bereitet. 
\emph{Noel Karlsson}
\index{Noel Karlsson}%
\index{Karlsson, Noel}%
und
\emph{Simon Widmer}
\index{Simon Widmer}%
\index{Widmer, Simon}%
zeigen in Kapitel~\ref{chapter:laplace}, wie die Laplace-Inversion
ausgehend von der Fourier-Umkehr auf das
\index{Fourier-Umkehr}%
Bromwich-Integral führt.
Die komplexe Analysis ermöglicht dann, aus diesem Integral eine effiziente
numerische Lösungsmethode zu finden, indem der Integrationspfad des
Bromwich-Integrals in die nach ihrem Erfinder die Talbot-Kontur genannt wird.
\index{Bromwich-Integral}%
\index{Talbot-Kontur}%

%
% weyl: Weyls Quadtree-Algorithmus
%
Die Bestimmung von Nullstellen eines Polynoms ist heute mit dem Computer
mit hoher Genauigkeit möglich.
Der newtonsche Algorithmus liefert beliebig viele Stellen einer Nullstelle
aber nur, wenn ein Startwert in der Nähe der Nullstelle gefunden werden kann.
Die ungefähre Lage der Nullstellen muss also vorgängig bestimmt werden.
\emph{Lukas Krüdewagen}
\index{Lukas Krüdewagen}%
\index{Krüdewagen, Lukas}%
erklärt in Kapitel~\ref{chapter:weyl}, wie moderne Algorithmen Ideen
verwenden, die auf Hermann Weyl zurückgehen und die komplexe Analysis
auf fundamentale Weise verwenden.


